\section{Introduction}

Heterogeneous catalysis, electrochemical energy conversion, and surface-mediated synthesis all depend critically on identifying the most stable adsorbate--surface configurations.
The configuration space for this problem grows combinatorially: site types (ontop, bridge, hollow, and their compositional variants), adsorbate orientation, and conformational flexibility combine to produce candidate counts that can exceed several hundred for a single intermetallic surface with a polyatomic adsorbate.
Traditional approaches---DFT enumeration, genetic algorithms, basin hopping---require many energy evaluations, each costing hundreds to thousands of CPU-hours, and scale poorly to the complex alloy and multimetallic surfaces that are of greatest practical interest.

\subsection{Machine learning surrogates and their limits}

Machine learning force fields (MLFFs) such as MACE~\cite{batatia2022mace}, CHGNet~\cite{deng2023chgnet}, and Equiformer~\cite{liao2023equiformerv2} have dramatically reduced the per-evaluation cost, achieving DFT-level energy predictions with orders-of-magnitude speedup.
However, MLFFs solve only the \textit{fast energy evaluation} problem; they do not address the upstream \textit{search strategy} problem of which configurations to evaluate in the first place.
Deep-learning-assisted configuration generators---graph neural networks that predict site stability, diffusion models that propose structures directly---remain open-loop: once trained, they cannot adapt based on new calculation results.
This architectural limitation causes systematic failures on surfaces where training-data coverage is sparse, precisely the complex compositions where better search is most needed.

\subsection{LLM-driven agents: opportunities and bottlenecks}

Large language models encode substantial chemical intuition from their training corpora and can generate plausible adsorption hypotheses from natural-language surface descriptions.
Adsorb-Agent~\cite{ock2024adsorbagent} demonstrated that an LLM agent can reduce the configuration search space by 27\% on a curated benchmark.
Yet Adsorb-Agent and similar systems operate in an \textit{open-loop} regime: the LLM proposes configurations, calculations execute, but the results are not fed back to the planner.
This means the agent cannot distinguish a thermodynamically unstable site from a calculation failure, cannot learn from dissociation or rearrangement events, and cannot correct the systematic biases that all current LLMs carry---preferring high-symmetry sites, underweighting surface composition effects, and ignoring adsorbate conformational flexibility.

\subsection{Closing the loop: physics-feedback-driven reasoning}

We introduce AdsMind, a closed-loop multi-agent framework in which each MLFF relaxation result feeds back into the LLM planner, enabling the system to detect its own reasoning errors, encode them as actionable search constraints, and decide autonomously when to stop.
Three mechanisms implement this feedback cycle:

\begin{enumerate}
    \item \textbf{Chemical Slip detection.} After each relaxation, the analyzer compares the planner's intended adsorption site with the realized coordination environment. A mismatch---termed a \textit{Chemical Slip}---signals that the potential energy surface (PES) topology disagrees with the LLM's chemical intuition at that point, providing a physically grounded diagnostic of planner error.
    \item \textbf{FORBID constraints.} Detected slips and dissociation events are converted into explicit negative constraints in the planner's prompt history, preventing the agent from retrying site families that have already been shown to be unstable or misclassified.
    \item \textbf{Autonomous termination.} The agent monitors whether successive iterations converge to the same best energy and terminates early when further exploration is unlikely to improve the result, avoiding unnecessary compute.
\end{enumerate}

We frame this design through a sequential-decision lens: adsorption configuration search is analogous to a partially observable Markov decision process (POMDP), where the PES is the hidden state, each relaxation is an observation, and the planner's next proposal is conditioned on the full history of prior observations.
Open-loop methods use a fixed policy $\pi(a \mid s)$ that cannot adapt after deployment; AdsMind approximates a history-dependent policy $\pi(a \mid h_t)$ through the LLM's in-context learning, without requiring explicit reward engineering or policy optimization.

\subsection{Hypotheses and evaluation}

We test two complementary claims:

\begin{description}
    \item[H$_1$ (search quality):] Full AdsMind, compared to the no-feedback single-shot baseline, finds lower-energy configurations ($\Delta E > 0.05$~eV) on a substantial fraction of intermetallic surfaces, with the benefit strongest when the base planner is least capable.
    \item[H$_2$ (backend robustness):] The iterative closed-loop architecture reduces energy variance across heterogeneous LLM backends, driving planners with different internal biases toward the same physical solution.
\end{description}

We evaluate these hypotheses on a 20-surface benchmark derived from Adsorb-Agent's test set, a controlled 5-case ablation study executed independently on three LLM backends (Gemini~2.5~Pro, Grok-4, and GPT-5.4), and DFT validation of representative configurations.
All experiments use a fixed MACE-small physical backend so that observed differences reflect search strategy, not energy-function choice.

AdsMind is complementary to, not competitive with, open-loop screening tools.
High-throughput methods are efficient for initial down-selection across thousands of candidate surfaces; closed-loop refinement adds value in the subsequent deep exploration of promising but poorly characterized systems---intermetallics, multimetallics, and novel compositions where LLM priors are least reliable.
